\documentclass[11pt]{article}
\usepackage[a4paper,margin=26mm]{geometry}
\usepackage{amsmath,amssymb,mathtools}
\usepackage{unicode-math}
\setmainfont{Latin Modern Roman}
\setmathfont{Latin Modern Math}
\usepackage{xurl,hyperref}
\hypersetup{hidelinks,pdftitle={Problem statement}}
\newcommand{\sref}[2]{\hyperlink{source:#1:#2}{[S#2]}}
\newcommand{\cataloguescope}[1]{#1}
\begin{document}
\section{Coarse foliations — Question 12.3}\label{statement}

\subsection{Definitions and mathematical statement}

\paragraph{Shared notation.}$\mathbb N=\{1,2,\ldots\}$, and $\mathbb Z,\mathbb Q,\mathbb R,\mathbb C$ denote the integers, rationals, reals and complex numbers. Any other conventions are specified locally.


A codimension-one foliation $\mathcal F$ decomposes a $3$-manifold $M$ into surface leaves, locally as the horizontal slices of a product chart. In the terminology of the source, it is taut if a closed curve transverse to $\mathcal F$ meets every leaf. A minimal foliation means that every leaf is dense; a minimal surface means a surface with vanishing mean curvature for the specified Riemannian metric. These two uses of minimal are distinct. A surface lamination $\Lambda$ is a closed subset of $M$ locally homeomorphic to a product of a disk with a closed transverse set. It is essential when it has no sphere leaf or torus leaf bounding a solid torus, and its complementary regions are irreducible, have incompressible boundary, and have no compressing monogon. Such a monogon is a compressing disk with one boundary arc on a boundary leaf and the other along a transverse interval fiber. It is genuine if some complementary region is not an interval bundle. A sublamination is a closed union of leaves.
\paragraph{Statement recorded in the supplied source.}
Develop a useful notion of inessential deformation for foliations and laminations. Determine which changes should preserve their essential geometric information, including monotone equivalence (collapsing foliated intervals), cutting and shearing along a surface or transverse lamination, moving a branch locus by isotopy in a branched cover, and isomorphism of universal circles. Formulate and justify appropriate equivalence relations instead of assuming all these operations automatically have identical effects. \sref{statement}{2}
\subsection{Short English statement}
Find the right equivalence relation for deformations of foliations and laminations, testing the source list of geometric and topological changes.
\subsection{Sources}
\begin{description}
\item[\hypertarget{source:statement:1}{S1}] ULAM.AI and the UnsolvedMath Contributors, \emph{UnsolvedMath: A Curated Collection of Open Mathematics Problems}, catalogue entry. CC BY 4.0. \url{https://huggingface.co/datasets/ulamai/UnsolvedMath}.
\item[\hypertarget{source:statement:2}{S2}] Danny Calegari, \emph{Problems in foliations and laminations of 3-manifolds}, version 0.78 (2002), Question 12.3; Definitions 1.1--1.10 and local definitions and remarks. \url{https://arxiv.org/abs/math/0209081v1}.
\end{description}
\label{end:statement}
\end{document}
